{\renewcommand{\Article}[1]{\par\medskip\Needspace{4\baselineskip}\pdfbookmark[1]{Art. #1}{erdmagnetismus-art-#1}\centerline{\textbf{#1}}\nobreak\smallskip}\Article{23}}
Es gibt aber ein anderes Verfahren, das von einem Teil dieser Schwierigkeiten frei ist und für den ersten Versuch besonders geeignet scheint. Wir wollen es hier entwickeln und dabei die Bedenken nennen, denen seine Anwendung unter den gegenwärtigen Umständen noch unterliegt. Dieses Verfahren setzt voraus, dass alle drei Elemente an Punkten bekannt sind, die auf einer hinreichenden Zahl von Breitenkreisen so verteilt sind, \SourceLeaf{151}dass jeder Breitenkreis dadurch in eine hinreichende Zahl gleicher Abschnitte geteilt wird.

Aus den in der üblichen Form gegebenen Elementen sind zuerst die Zahlenwerte von $X,Y$ und $Z$ abzuleiten.

Dann bringt man nach der bekannten Methode die Werte von $X,Y$ und $Z$ auf jedem Breitenkreis in die Form
\[
\begin{aligned}
X&=k+k'\cos\lambda+K'\sin\lambda+k''\cos2\lambda+K''\sin2\lambda\\&\quad+k'''\cos3\lambda+K'''\sin3\lambda+\text{u. s. w.}\\
Y&=l+l'\cos\lambda+L'\sin\lambda+l''\cos2\lambda+L''\sin2\lambda\\&\quad+l'''\cos3\lambda+L'''\sin3\lambda+\text{u. s. w.}\\
Z&=m+m'\cos\lambda+M'\sin\lambda+m''\cos2\lambda+M''\sin2\lambda\\&\quad+m'''\cos3\lambda+M'''\sin3\lambda+\text{u. s. w.}
\end{aligned}
\]
Man erhält also für jeden der Koeffizienten $k,l,m,k'$ usw. so viele Werte, wie Breitenkreise behandelt werden.

Der Theorie zufolge sollte auf jedem Breitenkreis $l=0$ gelten. Die aus der Rechnung hervorgehenden Werte von $l$ geben also bereits eine Art Maßstab dafür, wie unzuverlässig die zugrunde gelegten Zahlen noch sind.

Es gibt doppelt so viele Gleichungen
\[
\begin{aligned}
k&=-g^{1,0}\frac{\mathrm{d}P^{1,0}}{\mathrm{d}u}-g^{2,0}\frac{\mathrm{d}P^{2,0}}{\mathrm{d}u}-g^{3,0}\frac{\mathrm{d}P^{3,0}}{\mathrm{d}u}-\text{u. s. w.}\\
m&=2g^{1,0}P^{1,0}+3g^{2,0}P^{2,0}+4g^{3,0}P^{3,0}+\text{u. s. w.}
\end{aligned}
\]
wie Breitenkreise. Setzt man in $\frac{\mathrm{d}P^{1,0}}{\mathrm{d}u}$ usw. und in $P^{1,0}$ usw. die entsprechenden Zahlenwerte von $u$ ein, so bestimmt man daraus nach der Methode der kleinsten Quadrate von den Koeffizienten $g^{1,0},g^{2,0},g^{3,0}$ usw. so viele, wie berücksichtigt werden sollen.

Ebenso dienen die Gleichungen
\[
\begin{aligned}
-k'&=g^{1,1}\frac{\mathrm{d}P^{1,1}}{\mathrm{d}u}+g^{2,1}\frac{\mathrm{d}P^{2,1}}{\mathrm{d}u}+g^{3,1}\frac{\mathrm{d}P^{3,1}}{\mathrm{d}u}+\text{u. s. w.}\\
L'&=g^{1,1}\frac{P^{1,1}}{\sin u}+g^{2,1}\frac{P^{2,1}}{\sin u}+g^{3,1}\frac{P^{3,1}}{\sin u}+\text{u. s. w.}\\
m'&=2g^{1,1}P^{1,1}+3g^{2,1}P^{2,1}+4g^{3,1}P^{3,1}+\text{u. s. w.}
\end{aligned}
\]
deren Zahl insgesamt dreimal so groß ist wie die Zahl der Breitenkreise, zur Bestimmung der Koeffizienten $g^{1,1},g^{2,1},g^{3,1}$ usw. Die folgenden Gleichungen
\SourceLeaf{152}
\[
\begin{aligned}
-K'&=h^{1,1}\frac{\mathrm{d}P^{1,1}}{\mathrm{d}u}+h^{2,1}\frac{\mathrm{d}P^{2,1}}{\mathrm{d}u}+h^{3,1}\frac{\mathrm{d}P^{3,1}}{\mathrm{d}u}+\text{u. s. w.}\\
-l'&=h^{1,1}\frac{P^{1,1}}{\sin u}+h^{2,1}\frac{P^{2,1}}{\sin u}+h^{3,1}\frac{P^{3,1}}{\sin u}+\text{u. s. w.}\\
M'&=2h^{1,1}P^{1,1}+3h^{2,1}P^{2,1}+4h^{3,1}P^{3,1}+\text{u. s. w.}
\end{aligned}
\]
dienen entsprechend zur Bestimmung der Koeffizienten $h^{1,1}$, $h^{2,1}$, $h^{3,1}$ usw.

Ferner dienen zur Bestimmung der Koeffizienten $g^{2,2}$, $g^{3,2}$, $g^{4,2}$ usw. die Gleichungen
\[
\begin{aligned}
-k''&=g^{2,2}\frac{\mathrm{d}P^{2,2}}{\mathrm{d}u}+g^{3,2}\frac{\mathrm{d}P^{3,2}}{\mathrm{d}u}+g^{4,2}\frac{\mathrm{d}P^{4,2}}{\mathrm{d}u}+\text{u. s. w.}\\
L''&=2g^{2,2}\frac{P^{2,2}}{\sin u}+2g^{3,2}\frac{P^{3,2}}{\sin u}+2g^{4,2}\frac{P^{4,2}}{\sin u}+\text{u. s. w.}\\
m''&=3g^{2,2}P^{2,2}+4g^{3,2}P^{3,2}+5g^{4,2}P^{4,2}+\text{u. s. w.}
\end{aligned}
\]
und auf ähnliche Weise ergeben sich die Koeffizienten der folgenden höheren Ordnungen.

\Article{24}
Der Vorzug dieses Verfahrens gegenüber dem in Artikel 22 angegebenen besteht vor allem darin, dass die unbekannten Größen in Gruppen zerfallen, die jeweils für sich bestimmt werden. Das erleichtert die Rechnung außerordentlich, während beim anderen Verfahren die Vermischung sämtlicher Unbekannten ihre Trennung überaus mühsam macht. Dagegen hat dieses Verfahren den Nachteil, dass es keine unmittelbaren Beobachtungen zur Grundlage hat und seine Grundlagen grafischen Darstellungen entnehmen muss. Selbst dort, wo Beobachtungen vorliegen, können die Darstellungen sie nur grob wiedergeben. In Gegenden, in denen es weit und breit ganz an Beobachtungen fehlt, sind die Darstellungen lediglich aufgrund von Vermutungen und bis zu einem gewissen Grad willkürlich ergänzt und können daher sehr weit von der Wahrheit abweichen. Es bleibt jedoch nur die Wahl, entweder alle Versuche auszusetzen, bis wesentlich vollständigere und zuverlässigere Daten vorliegen, oder mit den derzeit noch so unsicheren Mitteln einen ersten Versuch zu wagen, von dem man wenig mehr als eine grobe Annäherung erwarten darf. Einen sicheren Maßstab für den Wert des Ergebnisses liefert im Nachhinein in jedem Fall der genaue Vergleich mit tatsächlichen Beobachtungen aus allen \SourceLeaf{153}Teilen der Erde. Ergibt diese Prüfung, dass der erste Versuch nicht völlig misslungen ist, so wird er eine große Hilfe dabei sein, künftige Versuche auf dem einen oder dem anderen Weg sinnvoll vorzubereiten.

\Article{25}
Schon vor vielen Jahren hatte ich wiederholt solche Versuche begonnen, musste sie aber immer wieder aufgeben, weil die verfügbaren Daten sich als viel zu spärlich erwiesen. Dennoch wäre ich schon früher geneigt gewesen, einen Versuch zu Ende zu führen, wenn mein mehrfach geäußerter Wunsch in Erfüllung gegangen wäre. dass die horizontalen Intensitäten für sich in einer allgemeinen Karte dargestellt würden, die sich durch eine Verbindung der vorhandenen unvollkommenen Übersichtskarten für die Inklination und die Gesamtintensität nicht ersetzen ließ.

Das Erscheinen der Karte der Gesamtintensität von \textsc{Sabine} (im siebten \textit{Report of the British association for the advancement of science}) hat mich nun veranlasst, einen neuen Versuch zu unternehmen und abzuschließen, der allerdings nur unter dem im vorangegangenen Artikel angegebenen Gesichtspunkt betrachtet werden soll.

Die der Rechnung zugrunde gelegten Daten wurden für die Intensität der genannten Karte entnommen, für die Deklination der Karte von \textsc{Barlow} (\textit{Philosophical Transactions} 1833) und für die Inklination der von \textsc{Horner} entworfenen Karte (Physikalisches Wörterbuch Band 6.), und zwar für jeweils zwölf Punkte auf sieben Breitenkreisen. Die Lücken, die diese Karten in weiten Gebieten lassen, konnten meist nur auf höchst unsichere Weise ergänzt werden.

Im Laufe der Rechnung zeigte sich bald, dass sie mindestens bis zu den Größen vierter Ordnung ausgedehnt werden musste, womit die Zahl der zu bestimmenden Koeffizienten auf 24 steigt. Aller Wahrscheinlichkeit nach werden auch die Glieder fünfter Ordnung noch ins Gewicht fallen. Bei einem ersten Versuch stehen die Werte von $k$, $m$, $k'$ usw. aber noch viel zu sehr unter dem Einfluss der vielen unzuverlässigen Daten, die ein solcher Versuch seiner Natur nach einschließen muss, als dass man eine noch größere Zahl von Unbekannten in die Elimination einbeziehen dürfte.

\SourceLeaf{154}
Es ist noch zu bemerken, dass die Intensitäten in \textsc{Sabine}s Karte in derselben willkürlichen Einheit angegeben sind, die bisher üblicherweise verwendet wurde und in der die Gesamtintensität in London $=1{,}372$ gesetzt wird. Diese Einheit ist hier bei der Berechnung der Koeffizienten ebenso wie in der weiter unten zu erklärenden Hilfstafel so abgeändert worden, dass alle Zahlen tausendmal so groß werden und damit die Intensität für London $=1372$ zugrunde liegt. Im Übrigen kann für die Intensität offenbar jede beliebige Einheit verwendet werden, da man auch die Einheit für $\mu$ als willkürlich ansehen und sie stets entsprechend wählen kann. Will man daraus weitere Folgerungen ziehen, für die $\mu$ in absolutem Maß ausgedrückt sein muss, braucht man nur sämtliche Koeffizienten mit demselben Faktor zu multiplizieren, der nötig ist, um die in jener Einheit ausgedrückten Intensitätszahlen in absolutes Maß umzurechnen.

\Article{26}
Bei der ersten Rechnung werden die Längen $\lambda$ von Greenwich nach Osten gezählt. Sie liefert für die 24 Koeffizienten die folgenden Zahlenwerte.
\[
\begin{array}{r|r}
g^{1,0}=+925{,}782&g^{2,2}=+0{,}493\\
g^{2,0}=-22{,}059&g^{3,2}=-73{,}193\\
g^{3,0}=-18{,}868&g^{4,2}=-45{,}791\\
g^{4,0}=-108{,}855&h^{2,2}=-39{,}010\\
g^{1,1}=+89{,}024&h^{3,2}=-22{,}766\\
g^{2,1}=-144{,}913&h^{4,2}=+42{,}573\\
g^{3,1}=+122{,}936&g^{3,3}=+1{,}396\\
g^{4,1}=-152{,}589&g^{4,3}=+19{,}774\\
h^{1,1}=-178{,}744&h^{3,3}=-18{,}750\\
h^{2,1}=-6{,}030&h^{4,3}=-0{,}178\\
h^{3,1}=+47{,}794&g^{4,4}=+4{,}127\\
h^{4,1}=+64{,}112&h^{4,4}=+3{,}175
\end{array}
\]
Diese Zahlen, die man als die \textit{Elemente der Theorie des Erdmagnetismus} ansehen kann, sind hier genau so eingesetzt und der später zu beschreibenden Hilfstafel zugrunde gelegt, wie die Rechnung sie ergeben hat, einschließlich der Dezimalbrüche. Für jeden, der mit solchen Rechnungen vertraut ist, erübrigt sich der \SourceLeaf{155}Hinweis, dass diese Bruchteile an sich keinen Wert haben, da wir noch weit davon entfernt sind, auch nur die ganzen Einer zuverlässig bestimmen zu können. Es ist aber wichtig, die Beobachtungen mit ein und demselben bestimmten System von Elementen genau zu vergleichen. Es bestand daher kein Grund, das Ergebnis der Rechnung zu ändern, da das Weglassen der Dezimalbrüche die Vergleichsrechnungen überhaupt nicht erleichtert hätte.

\Article{27}
Der nach den obigen Zahlen entwickelte Ausdruck für $V$ ist unten angegeben. Zur Abkürzung steht $e$ für $\cos u$ und $f$ für $\sin u$.
\[
\begin{aligned}
\frac VR={}&-1{,}977+937{,}103e+71{,}245ee-18{,}868e^3-108{,}855e^4\\
&+(64{,}437-79{,}518e+122{,}936ee+152{,}589e^3)f\cos\lambda\\
&+(-188{,}303-33{,}507e+47{,}794ee+64{,}112e^3)f\sin\lambda\\
&+(7{,}035-73{,}193e-45{,}791ee)ff\cos2\lambda\\
&+(-45{,}092-22{,}766e-42{,}573ee)ff\sin2\lambda\\
&+(1{,}396+19{,}774e)f^3\cos3\lambda\\
&+(-18{,}750-0{,}178e)f^3\sin3\lambda\\
&+4{,}127f^4\cos4\lambda\\
&+3{,}175f^4\sin4\lambda
\end{aligned}
\]

Hier sollen auch die vollständig entwickelten Ausdrücke für die drei Komponenten der magnetischen Kraft angegeben werden.
\[
\begin{aligned}
X={}&(937{,}103+142{,}490e-56{,}603ee-435{,}420e^3)f\\
&+(-79{,}518+181{,}435e-298{,}732ee\\
&\qquad-368{,}808e^3+610{,}357e^4)\cos\lambda\\
&+(-33{,}507+283{,}892e+259{,}349ee\\
&\qquad-143{,}383e^3-256{,}448e^4)\sin\lambda\\
&+(-73{,}193-105{,}652e+219{,}579ee+183{,}164e^3)f\cos2\lambda\\
&+(-22{,}766+175{,}330e+68{,}098ee-170{,}292e^3)f\sin2\lambda\\
&+(19{,}774-4{,}188e-79{,}096ee)ff\cos3\lambda\\
&+(-0{,}178+56{,}250e+0{,}716ee)ff\sin3\lambda\\
&-16{,}508ef^3\cos4\lambda\\
&-12{,}701ef^3\sin4\lambda
\end{aligned}
\]
\SourceLeaf{156}
\[
\begin{aligned}
Y={}&(188{,}303+33{,}507e-47{,}794ee-64{,}112e^3)\cos\lambda\\
&+(64{,}437-79{,}518e+122{,}936ee-152{,}589e^3)\sin\lambda\\
&+(90{,}184+45{,}532e-85{,}146ee)f\cos2\lambda\\
&+(14{,}070-146{,}386e-91{,}582ee)f\sin2\lambda\\
&+(56{,}250+0{,}534e)ff\cos3\lambda\\
&+(4{,}188+59{,}322e)ff\sin3\lambda\\
&-12{,}701f^3\cos4\lambda\\
&+16{,}508f^3\sin4\lambda
\end{aligned}
\]
\[
\begin{aligned}
Z={}&-24{,}593+1896{,}847e+400{,}343ee-75{,}471e^3-544{,}275e^4\\
&+(79{,}700-107{,}763e+491{,}744ee-762{,}946e^3)f\cos\lambda\\
&+(-395{,}724-155{,}473e+191{,}176ee+320{,}560e^3)f\sin\lambda\\
&+(34{,}187-292{,}772e-228{,}955ee)ff\cos2\lambda\\
&+(-147{,}439-91{,}064e+212{,}865ee)ff\sin2\lambda\\
&+(5{,}584+98{,}870e)f^3\cos3\lambda\\
&+(-75{,}000-0{,}890e)f^3\sin3\lambda\\
&+20{,}635f^4\cos4\lambda\\
&+15{,}876f^4\sin4\lambda
\end{aligned}
\]

Nachdem diese Komponenten für einen gegebenen Ort berechnet sind, erhält man die Bestimmungsgrößen der magnetischen Kraft in der üblichen Form auf folgende Weise. Es sei $\delta$ die Deklination, $i$ die Inklination, $\psi$ die Gesamtintensität und $\omega$ die horizontale Intensität. Man bestimmt zuerst $\delta$ und $\omega$ mithilfe der Formeln
\[
X=\omega\cos\delta,\qquad Y=\omega\sin\delta
\]
und dann $i$ und $\psi$ mithilfe der folgenden Formeln
\[
\omega=\psi\cos i,\qquad Z=\psi\sin i
\]

